\section*{Quadratic forms:}

Let A be a commutative ring. A quadratic form Q on an A-module E is a map from E into A such that \(\mathrm{Q}(\alpha x) = \alpha^{2}\mathrm{Q}(x)\) for \(\alpha \in A, x \in E\) , and that the map

\[
(x, y) \rightarrow \Phi (x, y) = Q (x + y) - Q (x) - Q (y)
\]

is a bilinear form (necessarily symmetric), said to be associated with the quadratic form Q. One has \(\Phi(x, x) = 2Q(x)\) ; conversely, for every bilinear form \(\Psi\) on E, \(x \to \Psi(x, x)\) is a quadratic form on E;

when A is a field of characteristic ≠ 2, quadratic forms and symmetric bilinear forms on E therefore correspond bijectively.

